\subsection{9.4: finite generators of group-isotope varieties}
\label{cand:9.4}
Use ordinary quasigroups in the three-operation, one-sorted signature
\((*,\backslash,/)\). If a nontrivial finite group generates a variety
\(\mathcal M\), the class \(I\mathcal M\) of all isotopes of groups in
\(\mathcal M\) cannot be generated by any finite quasigroup, or by any
finite family of finite quasigroups. The same obstruction applies to
quasivarieties and pseudovarieties.

Define terms \(t_0(x,y)=y\), \(t_{j+1}(x,y)=x*t_j(x,y)\).
Every quasigroup of size at most \(n\) satisfies
\[
 t_{n!}(x,y)=y,                                      \tag{***}
\]
because each left translation is a permutation whose cycle lengths
divide \(n!\). Identities persist under each of the three generation
operations just listed.

Choose a subgroup \(C_p\) of the nontrivial finite generating group.
For any \(N\ge1\), let \(m=N+1\), \(A=\F_p^m\), and let \(T\) cycle
its standard basis \(e_0,\ldots,e_{m-1}\).
The operation
\[
 a*b=a+T(b),\qquad
 a\backslash c=T^{-1}(c-a),\qquad c/b=c-T(b)
\]
defines a quasigroup isotopic to \(A\); the isotopy is
\((\mathrm{id},T,\mathrm{id})\). Only a subgroup and finite direct
powers of the original group were used, so this finite quasigroup
belongs to \(I\mathcal M\) in all three interpretations.
At \(x=0,y=e_0\), however,
\[
 t_N(0,e_0)=T^N(e_0)=e_N\ne e_0.
\]
Taking \(N=n!\) contradicts (***) for every proposed finite family of
finite generators.

There is also a local-finiteness obstruction. On
\(A=\bigoplus_{i\in\Z}\F_pe_i\), for the same prime \(p\) chosen above,
use the bilateral shift and the same
operation. This infinite quasigroup is generated by \(e_0\):
\(e_0\backslash e_0=0\), the terms \(0*e_i\) and
\(0\backslash e_i\) move in both directions, and
\(a+b=a*(0\backslash b)\) forms all finite sums.
A variety generated by a finite algebra \(Q\) is locally finite,
since its free algebra on \(k\) generators has at most
\(|Q|^{|Q|^k}\) term functions.

The interpretation is essential. In a three-sorted reversible automaton,
\(x*(x*y)\) need not be well typed. Gvaramiya distinguishes these
objects from ordinary quasigroups, and states the ordinary isotope
variety theorem in \cite[Section~3, Theorem~8]{gvaramiya-automata};
compare also \cite{krapez-marinkovic}.
The argument concerns ordinary algebraic generation, as in the supplied
statement, and makes no assertion about generation defined to include
closure under isotopy. Its simplicity makes independent checking of
the intended convention and priority particularly valuable.
See \artifact{research/9.4-proof.md}.
